\section{Introduction}

Ce projet a \'et\'e r\'ealis\'e dans le cadre du module NLP du Master~1. Il consiste \`a concevoir,
impl\'ementer et \'evaluer un syst\`eme de \emph{Retrieval-Augmented Generation} (RAG) : un assistant
capable de r\'epondre \`a des questions en langage naturel en s'appuyant sur un corpus documentaire
et un grand mod\`ele de langage (LLM).

\subsection{Contexte et objectifs}

Le domaine retenu est celui des \textbf{concepts fondamentaux du NLP et du machine learning}
(transformers, embeddings, recherche d'information, m\'etriques d'\'evaluation, etc.). Ce choix
permet de constituer un corpus coh\'erent \`a partir d'articles Wikipedia tout en restant directement
utile pour r\'eviser les notions vues dans le module.

L'objectif g\'en\'eral est de couvrir la cha\^ine compl\`ete d'un syst\`eme RAG :
\begin{itemize}[noitemsep]
  \item collecte et nettoyage des donn\'ees ;
  \item d\'ecoupage en \emph{chunks} et indexation vectorielle ;
  \item pipeline de r\'ecup\'eration et de g\'en\'eration ;
  \item interface utilisateur ;
  \item \'evaluation quantitative et visualisation des r\'esultats.
\end{itemize}

\subsection{P\'erim\`etre de ce rapport}

Ce rapport d\'ecrit l'\textbf{\'etat actuel} du projet au moment de la remise. Les phases
d'indexation, de pipeline RAG et d'\'evaluation automatique sont op\'erationnelles et produisent
des r\'esultats mesurables. L'interface Streamlit a \'egalement \'et\'e test\'ee et valid\'ee
(\texttt{streamlit run app.py}) ; seule la pr\'eparation de la d\'emonstration orale reste
\`a finaliser (voir section~\ref{sec:etat}). Le rapport est organis\'e comme suit :
\begin{enumerate}[noitemsep]
  \item \textbf{Architecture} du syst\`eme (section~\ref{sec:archi}) ;
  \item \textbf{Donn\'ees} et pr\'eparation du corpus (section~\ref{sec:donnees}) ;
  \item \textbf{Impl\'ementation} et choix techniques (section~\ref{sec:impl}) ;
  \item \textbf{M\'ethodologie d'\'evaluation} (section~\ref{sec:methode}) ;
  \item \textbf{R\'esultats} et graphiques (section~\ref{sec:resultats}) ;
  \item \textbf{Discussion} et \textbf{conclusion} (sections~\ref{sec:discussion}
        et~\ref{sec:conclusion}).
\end{enumerate}
